\documentclass[12pt]{article}
\usepackage[utf8]{inputenc}
\usepackage[T1]{fontenc}
\usepackage[a4paper,margin=1in]{geometry}
\usepackage{enumitem}
\usepackage{hyperref}
\usepackage{parskip}
\setlist[itemize]{leftmargin=*,itemsep=2pt,topsep=2pt}
\setlist[enumerate]{leftmargin=*,itemsep=2pt,topsep=2pt}
\begin{document}

\section*{GRYORK LSP Phasewise Implementation Plan}


\section{1) Purpose}


This document converts the full diagnostic in Gryork\_LSP\_Technical\_Diagnostic\_Report.md into an execution-grade, phasewise implementation plan for the LSP Compliance Transformation Programme.


Primary objective:

\begin{itemize}
\item Make Gryork NBFC-ready under RBI Digital Lending Guidelines (DLG), DPDPA-aligned consent architecture, and regulated audit/document standards.
\end{itemize}

Time horizon:

\begin{itemize}
\item Phase 1 (Weeks 1-12): Minimum NBFC-Ready State
\item Phase 2 (Weeks 13-28): Compliance Maturity
\item Phase 3 (Weeks 29-52): Scale Infrastructure
\end{itemize}

\section{2) Non-Negotiable Compliance Principles (Apply In All Phases)}


\begin{enumerate}
\item Compliance is architectural, not manual.
\item No disbursement without CEE gate pass.
\item No data sharing without purpose-specific valid consent.
\item No legal reliance on mutable logs.
\item NBFC operations must remain possible even if Gryork core systems are degraded.
\item Borrower rights (KFS disclosure, cooling-off, exit) must be system-enforced.
\item All API failures must enter auditable fallback states.
\item Separation of borrower, NBFC, and Gryork data/identity/access domains is mandatory.
\end{enumerate}

\section{3) Coverage Matrix (Source Section -> Plan Coverage)}


\begin{itemize}
\item Section 1 (Executive Diagnostic): Program baseline, risk mitigation sequencing, readiness gates.
\item Section 2 (LSP Requirements): Target controls embedded in architecture and backlog.
\item Section 3 (Gap A-H): Each gap mapped to mandatory implementation tracks.
\item Section 4 (Target Architecture): Multi-interface + GRYLINK Core + API Gateway design implemented in phases.
\item Section 5 (Core Modules): Module-by-module engineering backlog and acceptance criteria.
\item Section 6 (API Integration Layer): Integration catalog, trigger points, standards, resilience, and test requirements.
\item Section 7 (CEE): Hard-gate enforcement for DLG controls, direct disbursement, fee separation, consent checks.
\item Section 8 (IDR + Platform Failure): Independent WORM/TSA repository and 30-minute NBFC access SLA.
\item Section 9 (Audit): Hash-chain + WORM + TSA architecture, event completeness, retention requirements.
\item Section 10 (Escrow): Real-time monitoring, reconciliation, mismatch alerts, NBFC visibility.
\item Section 11 (Security): RBAC, data segregation, encryption, mTLS, key rotation, consent-gated data flow.
\item Section 12 (CTO Responsibility): certification checkpoints and audit-readiness governance.
\item Section 13 (Roadmap): Phase-by-phase implementation schedule aligned to diagnostic.
\item Section 14 (Team Instructions): backend/frontend/devops/product execution rules and quality bars.
\end{itemize}

\section{4) Gap-to-Workstream Mapping (A-H)}


\begin{itemize}
\item Gap A (No Compliance Enforcement Engine) -> Workstream WS-CEE
\item Gap B (Non-Immutable Audit Log) -> Workstream WS-AUDIT
\item Gap C (No Independent Document Repository) -> Workstream WS-IDR
\item Gap D (Absent Consent Management) -> Workstream WS-CONSENT
\item Gap E (No DLG Workflow Controls) -> Workstream WS-DLG
\item Gap F (Weak Escrow Monitoring) -> Workstream WS-ESCROW
\item Gap G (No Independent NBFC Access Layer) -> Workstream WS-NBFC
\item Gap H (Incomplete API Integration Architecture) -> Workstream WS-API
\end{itemize}

\section{5) Program Structure}


\subsection{5.1 Workstreams}


\begin{itemize}
\item WS-ARCH: Target architecture and service boundaries.
\item WS-CEE: Compliance Enforcement Engine and hard gates.
\item WS-CONSENT: Consent capture, storage, revocation, validation.
\item WS-KFS: KFS generation, delivery proof, cooling-off orchestration.
\item WS-IDR: Independent Document Repository with WORM + TSA.
\item WS-AUDIT: Immutable audit engine and evidence exports.
\item WS-API: External integrations and fallback/error architecture.
\item WS-KYC: CKYC/Aadhaar/PAN/MCA21 and identity controls.
\item WS-GST: GSTN validation and duplicate invoice checks.
\item WS-EPC: EPC confirmation with digital signatures.
\item WS-DOC: OCR + document validation + fraud detection.
\item WS-ESCROW: Escrow/NPCI monitoring and reconciliation.
\item WS-NBFC: Independent NBFC dashboard and APIs.
\item WS-SEC: RBAC, segregation, encryption, mTLS, key/secrets lifecycle.
\item WS-SRE: Reliability, observability, DR/BCP, SLA operations.
\item WS-PRODUCT: Regulatory communication, UX compliance, RIA governance.
\end{itemize}

\subsection{5.2 Milestone Gates}


\begin{itemize}
\item Gate G0 (Week 2): Architecture and control blueprint approved.
\item Gate G1 (Week 12): Minimum NBFC-ready controls live and certified.
\item Gate G2 (Week 28): Compliance maturity stack complete and walkthrough passed.
\item Gate G3 (Week 52): Scale, security certification, inspection readiness achieved.
\end{itemize}

\section{6) Phase 1 - Minimum NBFC-Ready State (Weeks 1-12)}


\section{6.1 Phase Goal}


Deliver the minimum enforceable compliance architecture before onboarding any NBFC.


\section{6.2 Mandatory Outcomes (From Diagnostic)}


\begin{enumerate}
\item Independent Document Repository live (WORM + TSA + NBFC direct credentials).
\item Consent Management System live (purpose-specific capture and revocation).
\item KFS generation and delivery proof live in RBI-compliant format.
\item Cooling-off enforcement as hard workflow stop.
\item Bank account verification integrated for disbursement routing.
\item Audit Trail Engine live with hash-chain and first TSA anchors.
\item Independent NBFC Access Dashboard live.
\item DLG direct-disbursement routing validation enforced.
\item CKYC/PAN/Aadhaar integrations live in KYC module.
\item CEE gate checks for KFS/cooling-off/consent live.
\item Internal compliance walkthrough and CTO sign-off completed.
\end{enumerate}

\section{6.3 Detailed Backlog (By Sprint)}


\subsection{Sprint 1 (Weeks 1-2): Foundation and control architecture}


\begin{itemize}
\item P1-ARCH-01: Define service boundaries for GRYLINK Core services.
\item P1-ARCH-02: Finalize event taxonomy for end-to-end audit logging.
\item P1-ARCH-03: Define authoritative state machine and transition contract.
\item P1-SEC-01: Define identity domains (Borrower/NBFC/Gryork/Admin).
\item P1-SEC-02: Secrets vault baseline, certificate policy, key rotation schedules.
\item P1-SRE-01: Environments separation policy (dev/stage/prod), deployment approvals.
\item P1-PROD-01: Regulatory communication templates and content governance model.
\end{itemize}

Acceptance criteria:

\begin{itemize}
\item Architecture decision records approved.
\item Transition catalog includes all regulated steps.
\item No direct state mutation path bypassing orchestrator.
\end{itemize}

\subsection{Sprint 2 (Weeks 3-4): CEE + Consent + KFS foundations}


\begin{itemize}
\item P1-CEE-01: Implement CEE as mandatory pre-transition gate service.
\item P1-CEE-02: Enforce consent validity check before any cross-context data transfer.
\item P1-CEE-03: Implement direct-disbursement account validation gate.
\item P1-CEE-04: Implement fee-separation disbursement amount check.
\item P1-CMS-01: Purpose-specific consent schema with timestamp, IP, device/session.
\item P1-CMS-02: Consent revocation API with immediate effect.
\item P1-KFS-01: KFS template engine per RBI-required fields.
\item P1-KFS-02: Delivery proof ingestion (email/SMS/in-app acknowledgement).
\end{itemize}

Acceptance criteria:

\begin{itemize}
\item Disbursement cannot proceed when consent/KFS/cooling-off gates fail.
\item Consent revocation instantly blocks future sharing for that purpose.
\item KFS payload completeness validation implemented.
\end{itemize}

\subsection{Sprint 3 (Weeks 5-6): IDR + Audit core}


\begin{itemize}
\item P1-IDR-01: Stand up independent IDR infra domain, auth, metadata store.
\item P1-IDR-02: Enable WORM policy and immutable retention controls.
\item P1-IDR-03: Implement document SHA-256 fingerprinting pipeline.
\item P1-IDR-04: Integrate TSA RFC 3161 timestamp token generation at upload.
\item P1-IDR-05: NBFC credentialed read APIs independent of Gryork auth.
\item P1-IDR-06: Enforce minimum 8-year retention policy with lifecycle controls.
\item P1-IDR-07: Implement indexed retrieval and delivery workflow to meet <=30-minute NBFC access SLA.
\item P1-AUD-01: Hash-chain append-only event store.
\item P1-AUD-02: TSA anchoring job for audit batches.
\item P1-AUD-03: Event integrity verifier tool (chain + TSA checks).
\item P1-AUD-04: Implement retention/search/archive policy enabling <=48-hour retrieval across retention window.
\end{itemize}

Acceptance criteria:

\begin{itemize}
\item Document modification/delete attempts are blocked and audited.
\item NBFC can retrieve documents directly without Gryork admin involvement.
\item Audit chain verification passes end-to-end test.
\item IDR retention, retrieval, and access SLAs validated in controlled test runs.
\end{itemize}

\subsection{Sprint 4 (Weeks 7-8): KYC and mandatory API integrations}


\begin{itemize}
\item P1-KYC-01: CKYC integration with consent checks and refresh logic.
\item P1-KYC-02: Aadhaar authentication integration (masked storage only).
\item P1-KYC-03: PAN verification integration and mismatch handling.
\item P1-KYC-04: KYC signed result record pushed to IDR.
\item P1-API-01: Bank account verification integration and verified-account registry.
\item P1-API-02: API circuit breaker, timeout/retry, fallback state framework.
\item P1-API-03: API request/response metadata logging to audit engine.
\end{itemize}

Acceptance criteria:

\begin{itemize}
\item Raw PII responses not persisted in plain logs.
\item All integration failures create auditable fallback states.
\item Verified account registry used by CEE disbursement gate.
\end{itemize}

\subsection{Sprint 5 (Weeks 9-10): NBFC independent access and DLG enforcement}


\begin{itemize}
\item P1-NBFC-01: Launch independent NBFC dashboard auth and role model.
\item P1-NBFC-02: Portfolio view, document fetch, event notifications.
\item P1-DLG-01: Cooling-off countdown and immutable start trigger from KFS confirmation.
\item P1-DLG-02: Exit request flow that hard-cancels disbursement path.
\item P1-DLG-03: Disbursement routing proof record and signed compliance certificate.
\item P1-DLG-04: CEE enforcement for KFS term mismatch vs disbursement payload.
\end{itemize}

Acceptance criteria:

\begin{itemize}
\item NBFC can perform core review actions without Gryork operator support.
\item Any non-compliant disbursement instruction is blocked.
\item Exit during cooling-off immediately transitions to cancellation state.
\end{itemize}

\subsection{Sprint 6 (Weeks 11-12): Hardening, walkthrough, certification}


\begin{itemize}
\item P1-QA-01: Compliance regression suite for all CEE hard gates.
\item P1-QA-02: Audit completeness validation against event taxonomy.
\item P1-QA-03: IDR access SLA drill and failure-mode tests.
\item P1-CTO-01: Internal compliance walkthrough with evidence pack.
\item P1-CTO-02: CTO readiness certification for minimum NBFC onboarding state.
\end{itemize}

Evidence required at Gate G1:

\begin{itemize}
\item CEE pass/fail logs for controlled test scenarios.
\item KFS delivery evidence and cooling-off enforcement logs.
\item Consent lifecycle evidence (capture/revoke/deny).
\item IDR independent access demonstration logs.
\item Audit chain integrity verification report.
\end{itemize}

\section{7) Phase 2 - Compliance Maturity (Weeks 13-28)}


\section{7.1 Phase Goal}


Complete full regulated operating capabilities required for scaled NBFC due diligence.


\section{7.2 Mandatory Outcomes (From Diagnostic)}


\begin{enumerate}
\item Account Aggregator consent-gated integration.
\item Credit bureau integrations with consent gate enforcement.
\item GSTN validation with duplicate detection.
\item EPC confirmation engine with digital signatures.
\item CERSAI search/registration integration.
\item Document Validation Engine (OCR + cross-reference).
\item Fraud Rule Engine with escalation workflows.
\item MCA21 entity verification integration.
\item Escrow Monitoring Engine with real-time reconciliation.
\item Field-level PII encryption and mTLS rollout.
\item Quarterly internal audit review process operational.
\item First full NBFC compliance walkthrough completed.
\end{enumerate}

\section{7.3 Detailed Backlog (By Capability Track)}


\subsection{Track A: Integration completion and reliability}


\begin{itemize}
\item P2-API-01: Implement Account Aggregator data pull and consent linkage.
\item P2-API-02: Integrate CIBIL/Equifax/CRIF with stage-gated triggers.
\item P2-API-03: Integrate GSTN APIs and invoice cross-match checks.
\item P2-API-04: Integrate MCA21 entity/charge lookup and refresh policy.
\item P2-API-05: Integrate CERSAI pre-decision search and post-sanction registration.
\item P2-API-06: Expand e-sign/DSC/eStamp service orchestration.
\item P2-API-07: Standardize integration telemetry, error classes, and operational dashboards.
\end{itemize}

Acceptance criteria:

\begin{itemize}
\item Each API has success, timeout, malformed, and partial-data test coverage.
\item All API output artifacts are traceable in audit + IDR references.
\end{itemize}

\subsection{Track B: Document, fraud, and underwriting quality}


\begin{itemize}
\item P2-DOC-01: OCR extraction pipeline for invoices/certificates/statements.
\item P2-DOC-02: Data consistency rules (GSTIN, dates, amounts, entity links).
\item P2-DOC-03: Fraud Rules Engine with anomaly scoring and escalation queues.
\item P2-DOC-04: Duplicate invoice detection within Gryork platform context.
\item P2-EPC-01: EPC confirmation workflow with legally valid digital signatures.
\end{itemize}

Acceptance criteria:

\begin{itemize}
\item Document acceptance requires validation report artifact.
\item Fraud flags generate traceable workflow states and escalation events.
\end{itemize}

\subsection{Track C: Escrow integrity and NBFC transparency}


\begin{itemize}
\item P2-ESC-01: Integrate escrow/bank/NPCI webhooks for debit/credit events.
\item P2-ESC-02: Matching engine vs expected transaction schedule.
\item P2-ESC-03: Unmatched/late transaction alerts to NBFC and Gryork ops.
\item P2-ESC-04: Daily reconciliation report generation and IDR archival.
\item P2-NBFC-01: NBFC dashboard expansion for escrow views and compliance reports.
\end{itemize}

Acceptance criteria:

\begin{itemize}
\item NBFC can independently download daily reconciliations from IDR.
\item Alerting SLOs defined and monitored for transaction mismatch events.
\end{itemize}

\subsection{Track D: Security and governance maturity}


\begin{itemize}
\item P2-SEC-01: Field-level encryption rollout for Aadhaar/PAN/account/mobile PII.
\item P2-SEC-02: mTLS between internal microservices.
\item P2-SEC-03: RBAC hardening with least privilege and privileged action approvals.
\item P2-SEC-04: Consent gate enforced in all cross-context data access paths.
\item P2-GOV-01: Quarterly internal audit process (engineering + compliance + ops).
\end{itemize}

Acceptance criteria:

\begin{itemize}
\item No unencrypted PII fields in primary stores.
\item Inter-service traffic without mTLS is blocked.
\item Privilege escalation is dual-authorized and fully audited.
\end{itemize}

\subsection{Track E: Controlled readiness and assurance}


\begin{itemize}
\item P2-QA-01: End-to-end compliance test harness across full loan lifecycle.
\item P2-QA-02: Failure-injection tests for API/provider outages.
\item P2-QA-03: First NBFC compliance walkthrough with evidence closure.
\end{itemize}

Evidence required at Gate G2:

\begin{itemize}
\item API integration conformance matrix and failover test results.
\item Escrow reconciliation completeness metrics.
\item Security control verification (field encryption + mTLS + RBAC).
\item Internal audit report and closure logs.
\end{itemize}

\section{8) Phase 3 - Scale Infrastructure (Weeks 29-52)}


\section{8.1 Phase Goal}


Scale from compliant platform to inspection-ready, high-resilience, multi-NBFC infrastructure.


\section{8.2 Mandatory Outcomes (From Diagnostic)}


\begin{enumerate}
\item Video KYC for high-value/non-OTP scenarios.
\item Multi-NBFC portfolio segregation and controls.
\item Advanced ML fraud detection with retraining pipeline.
\item Automated RBI-format regulatory reporting.
\item Cross-region DR for IDR and Audit with tested recovery.
\item ISO 27001 certification completion.
\item Rate limiting, DDoS protection, and WAF on external endpoints.
\item Borrower consent dashboard with full history and revocation.
\item Portfolio-level risk analytics for NBFC credit committees.
\item Bug bounty and external penetration testing.
\item Certified BCP/DR procedures.
\item RBI inspection readiness evidence pack complete.
\end{enumerate}

\section{8.3 Detailed Backlog (By Quarter)}


\subsection{Q3 (Weeks 29-40): Scale capability build-out}


\begin{itemize}
\item P3-KYC-01: Video KYC orchestration and evidence retention.
\item P3-NBFC-01: Tenant-aware, fully segregated multi-NBFC architecture.
\item P3-DOC-01: ML fraud model pipeline with drift detection/retraining.
\item P3-REP-01: Automated regulatory report generator (RBI-aligned formats).
\item P3-SEC-01: Edge security stack (WAF, DDoS controls, API rate limits).
\item P3-CONSENT-01: Borrower self-service consent dashboard.
\item P3-ANALYTICS-01: NBFC portfolio risk insights and committee views.
\end{itemize}

Acceptance criteria:

\begin{itemize}
\item Multi-tenant segregation tests pass for data and access boundaries.
\item Regulatory reports are reproducible and audit-linked.
\item Fraud model retraining events are logged and reviewable.
\end{itemize}

\subsection{Q4 (Weeks 41-52): Resilience, certification, and inspection readiness}


\begin{itemize}
\item P3-SRE-01: Cross-region replication for IDR and audit stores.
\item P3-SRE-02: DR drills with RTO/RPO measurement and remediation.
\item P3-GRC-01: ISO 27001 control implementation and certification audit.
\item P3-GRC-02: Bug bounty launch and external penetration test remediation.
\item P3-GRC-03: BCP/DR documentation finalization and sign-offs.
\item P3-CTO-01: RBI inspection evidence pack assembly and dry-run review.
\end{itemize}

Acceptance criteria:

\begin{itemize}
\item DR test results meet target RTO/RPO commitments.
\item External assessment findings closed or risk-accepted via governance.
\item Inspection package includes traceable evidence for all mandatory controls.
\end{itemize}

Evidence required at Gate G3:

\begin{itemize}
\item ISO 27001 certification artifacts.
\item DR and failover evidence.
\item Pen-test and bug bounty closure report.
\item End-to-end compliance evidence index for regulator/NBFC review.
\end{itemize}

\section{9) External API Integration Plan (From Section 6)}


\section{9.1 API Catalog and Trigger Points}


\begin{itemize}
\item GST/GSTN API: invoice submission and pre-committee re-validation.
\item CKYC Registry: onboarding and 12-month refresh.
\item UIDAI Aadhaar Auth: individual KYC and threshold-based requirements.
\item NSDL PAN Verify: onboarding, pre-assessment, pre-execution checks.
\item Account Aggregator: post-credit-assessment consented data pull.
\item Credit Bureau (CIBIL/Equifax/CRIF): post-KYC and consented pre-sanction checks.
\item MCA21: entity onboarding and pre-disbursement refresh.
\item CERSAI: pre-decision search and post-execution registration.
\item E-Sign/DSC/eStamp: KFS/loan/consent execution events.
\item TSA (RFC 3161): document upload and audit batch anchoring.
\item Bank Account Verify: onboarding and mandatory disbursement routing validation.
\item Escrow/NPCI APIs: pre-disbursement and continuous reconciliation events.
\item OCR/Document AI: document ingestion and lifecycle validation.
\item Fraud Detection Engine: KYC/doc/data inconsistency events.
\end{itemize}

\section{9.2 Integration Engineering Standards (Mandatory)}


\begin{itemize}
\item Credentials only in secrets vault; never in source code or committed env files.
\item Certificate-based auth where possible; shared secrets rotate <= 90 days.
\item Circuit breakers, retries, timeout budgets, and explicit fallback states.
\item Full audit records for request metadata, response status, trigger context, and timestamps.
\item PII-safe logging and encrypted partition storage.
\item Test suites: success, provider error, timeout, malformed response, partial-data handling.
\end{itemize}

\section{10) Team-Wise Responsibilities (From Section 14)}


\section{10.1 Backend Engineering}


\begin{itemize}
\item Enforce authN/authZ/input validation/audit logging/error discipline on all endpoints.
\item Keep workflow orchestrator as sole state-transition authority.
\item Ensure no silent failures on external API errors.
\item Implement schema-level data sovereignty and field-level PII encryption.
\end{itemize}

\section{10.2 Frontend/Mobile}


\begin{itemize}
\item Build plain-language consent UX with symmetric accept/reject paths.
\item Ensure KFS disclosure readability and delivery acknowledgement integrity.
\item Display cooling-off countdown and accessible exit actions.
\item Emit frontend action events for critical NBFC and borrower actions.
\end{itemize}

\section{10.3 DevOps/Infrastructure}


\begin{itemize}
\item Enforce strict environment isolation and production deployment approvals.
\item Deploy dedicated infra domains for IDR and Audit engines.
\item Maintain 99.99\% reliability posture for compliance-critical components.
\item Operate secrets management, monitoring, alerting, and DR drills.
\end{itemize}

\section{10.4 Product}


\begin{itemize}
\item Run Regulatory Impact Assessment for every feature.
\item Prioritize compliance controls above UX enhancements when conflict exists.
\item Co-own regulated borrower communications with compliance sign-off.
\item Treat KFS/consent/cooling-off messages as controlled regulatory artifacts.
\end{itemize}

\section{11) Definition of Done (DoD) for Regulated Features}


A feature is complete only when all are true:


\begin{enumerate}
\item Functional behavior implemented and validated.
\item Compliance gate logic encoded and passing required checks.
\item Audit events emitted with required identity/session/context fields.
\item Artifacts persisted to IDR/audit references where required.
\item Security controls (RBAC, encryption, consent gate) verified.
\item Failure modes tested with auditable fallback behavior.
\item Monitoring and alerting configured with ownership.
\item Compliance and product sign-off captured for regulated user communication changes.
\end{enumerate}

\section{12) Program Risk Register and Mitigations}


\begin{itemize}
\item Risk: Delayed external API onboarding.
\begin{itemize}
\item Mitigation: Sandbox-first integration, provider SLAs, fallback manual-review queues with audit coverage.
\end{itemize}
\end{itemize}

\begin{itemize}
\item Risk: Partial compliance implementation mistaken as ready state.
\begin{itemize}
\item Mitigation: Hard gate checklists for G1/G2/G3 and CTO written certifications.
\end{itemize}
\end{itemize}

\begin{itemize}
\item Risk: Data leakage via logs or cross-context queries.
\begin{itemize}
\item Mitigation: PII redaction pipeline, field-level encryption, consent-guarded data access middleware.
\end{itemize}
\end{itemize}

\begin{itemize}
\item Risk: IDR/Audit availability gaps undermine trust.
\begin{itemize}
\item Mitigation: Dedicated infra, independent auth, cross-region replication, regular failover tests.
\end{itemize}
\end{itemize}

\begin{itemize}
\item Risk: Team drift from regulated build practices.
\begin{itemize}
\item Mitigation: Section-14-aligned engineering standards baked into PR checks, release gates, and architecture review.
\end{itemize}
\end{itemize}

\section{13) Readiness Gate Checklists}


\section{13.1 Gate G1 (Week 12) - NBFC Onboarding Eligibility}


Must be live and validated:

\begin{itemize}
\item CEE
\item CMS
\item KFS + Cooling-off enforcement
\item IDR with NBFC direct access
\item Audit hash-chain + TSA anchors
\item Independent NBFC dashboard
\item Critical APIs: KYC, bank verification, e-sign
\item CTO certification and internal walkthrough evidence
\end{itemize}

\section{13.2 Gate G2 (Week 28) - Compliance Maturity}


Must be live and validated:

\begin{itemize}
\item AA, bureau, GSTN, MCA21, CERSAI
\item Document validation + fraud rules
\item Escrow monitoring + daily reconciliation
\item Field-level encryption + mTLS + RBAC hardening
\item Quarterly internal audit operation
\item First full NBFC walkthrough closure
\end{itemize}

\section{13.3 Gate G3 (Week 52) - Scale and Inspection Readiness}


Must be live and validated:

\begin{itemize}
\item Video KYC
\item Multi-NBFC segregation
\item Advanced ML fraud and retraining governance
\item Automated RBI-format reporting
\item DR/BCP and cross-region resilience
\item ISO 27001 certification
\item WAF/DDoS/rate limiting controls
\item Pen-test and bug bounty closure
\item Full regulator-ready evidence pack
\end{itemize}

\section{14) Immediate Execution Plan (Next 14 Days)}


\begin{enumerate}
\item Approve architecture and control blueprint (WS-ARCH, WS-SEC, WS-CEE).
\item Freeze authoritative workflow state machine and transition gate contract.
\item Create compliance event taxonomy and evidence index schema.
\item Stand up IDR and audit infrastructure foundations in staging.
\item Implement first CEE gate bundle (consent + KFS + cooling-off placeholders).
\item Start KYC integration adapters (CKYC/PAN/Aadhaar) with audit wrappers.
\item Define CTO readiness dashboard with Gate G1 criteria status.
\end{enumerate}

\section{15) Final Instruction}


If any design choice creates ambiguity between speed and compliance, choose the stricter compliance path and escalate to compliance review before implementation.

\end{document}


